\documentclass[14pt]{extarticle}
\input{../../../scripts/tex_preamble.tex}

\title{نظام إدارة النادي الرياضي - البرمجة الموجهة للكائنات (OOP)}

\begin{document}

\maketitle
\thispagestyle{fancy}

\begin{enumerate}[itemsep=3em]
    \item
    \textbf{بناء فئة الاشتراك (\textenglish{Member}):} \\
    اكتب الفئة \textenglish{Member} بحيث تشمل الخصائص التالية:
    \begin{itemize}
        \item \textenglish{MemberName} (اسم المشترك - نص).
        \item \textenglish{SubscriptionId} (رقم الاشتراك - عدد صحيح).
        \item \textenglish{IsActive} (حالة الفعالية - قيمة منطقية \textenglish{bool}).
    \end{itemize}
    أضف دالة بناء (Constructor) تستقبل القيم وتهيئ الخصائص المذكورة أعلاه.

    \item
    \textbf{بناء فئة البوابة (\textenglish{ClubGate}):} \\
    اكتب الفئة \textenglish{ClubGate} بحيث تشمل الخصائص التالية:
    \begin{itemize}
        \item \textenglish{GateName} (اسم البوابة - نص).
        \item \textenglish{Members} (مصفوفة من نوع \textenglish{Member} تمثل قائمة المشتركين).
    \end{itemize}
    أضف دالة بناء (Constructor) لتهيئة اسم البوابة والمصفوفة.

    أضف الدالة \textenglish{CheckEntry(int subId)} التي تستقبل رقم اشتراك، تبحث عنه في المصفوفة، وتفحص خاصية \textenglish{IsActive} الخاصة به. إذا كانت \textenglish{true} تطبع "تم فتح البوابة للمشترك [الاسم]"، وإذا كانت \textenglish{false} (أو كان الاشتراك غير موجود) تطبع "عذراً، الدخول غير مسموح". \\

    أضف دالة \textenglish{CheckEntry(Member member)} كتحميل زائد (Overload) تقوم بنفس العمل ولكنها تستقبل كائن المشترك مباشرة بدلاً من البحث عنه بواسطة رقم الاشتراك.

    \item
    \textbf{فحص الكود (في \textenglish{Main}):} \\
    اكتب برنامجًا رئيسيًا يقوم بما يلي:
    \begin{itemize}
        \item إنشاء بوابة (\textenglish{Gate A}) وتهيئة مصفوفة المشتركين بداخلها.
        \item إنشاء اشتراك فعّال باسمك، واشتراك غير فعّال باسم صديقك، وإضافتهما للمصفوفة.
        \item استدعاء دالة \textenglish{CheckEntry} وتمرير رقم الاشتراك الخاص بك لفحص الدخول.
        \item استدعاء دالة \textenglish{CheckEntry} وتمرير كائن الاشتراك الخاص بصديقك (باستخدام الـ Overload) والتأكد من المخرجات.
    \end{itemize}
\end{enumerate}

\clearpage
\section*{المرحلة الثانية: بناء فئات الموظفين (تكرار الكود)}
الهدف: إضافة أقسام الإدارة والتدريب للنظام باستخدام ما تعلمناه حتى الآن.

\begin{enumerate}[itemsep=3em]
    \item
    \textbf{بناء فئة المدرب (\textenglish{Trainer}):} \\
    اكتب الفئة \textenglish{Trainer} بحيث تشمل الخصائص التالية:
    \textenglish{Id} (رقم الهوية)، \textenglish{FullName} (الاسم الكامل)، \textenglish{PhoneNumber} (رقم الهاتف)، \textenglish{BasicSalary} (الراتب الأساسي)، \textenglish{Specialty} (التخصص). \\
    أضف المهام التالية:
    \begin{itemize}
        \item دالة بناء تستقبل 5 معطيات لتهيئة جميع الخصائص.
        \item دالة \textenglish{PrintInfo()} لطباعة بيانات المدرب.
        \item دالة \textenglish{CalculateSalary()} تُرجع قيمة الراتب الأساسي.
    \end{itemize}

    \item
    \textbf{بناء فئة المدير (\textenglish{Manager}):} \\
    اكتب الفئة \textenglish{Manager} بحيث تشمل الخصائص التالية:
    \textenglish{Id}، \textenglish{FullName}، \textenglish{PhoneNumber}، \textenglish{BasicSalary}، \textenglish{Department} (القسم). \\
    أضف المهام التالية:
    \begin{itemize}
        \item دالة بناء تستقبل 5 معطيات لتهيئة الخصائص.
        \item دالة \textenglish{PrintInfo()} لطباعة بيانات المدير.
        \item دالة \textenglish{CalculateSalary()} تُرجع قيمة الراتب الأساسي.
    \end{itemize}

    \item
    \textbf{فحص الكود (في \textenglish{Main}):} \\
    أضف للبرنامج الرئيسي الكود التالي:
    \begin{itemize}
        \item إنشاء كائن من نوع \textenglish{Trainer} وكائن من نوع \textenglish{Manager}.
        \item استدعاء \textenglish{PrintInfo()} و \textenglish{CalculateSalary()} لكليهما.
    \end{itemize}
\end{enumerate}

\clearpage
\section*{المرحلة الثالثة: صدمة التعديل المفاجئ}
أرسلت إدارة النادي طلبًا عاجلًا لتحديث النظام بالمتطلبات التالية:
\begin{itemize}
    \item إضافة خاصية إجبارية جديدة وهي البريد الإلكتروني (\textenglish{EmailAddress}) لجميع العاملين.
    \item إضافة علاوة بنسبة 10\% على الراتب لجميع العاملين بمناسبة الأعياد.
\end{itemize}

\begin{enumerate}[itemsep=3em]
    \item
    \textbf{المهام (العمل اليدوي المكرر):} \\
    قم بتحديث الفئات السابقة كالتالي:
    \begin{itemize}
        \item افتح فئة \textenglish{Trainer} وأضف الخاصية الجديدة (\textenglish{EmailAddress}).
        \item عدّل دالة البناء في \textenglish{Trainer} لتستقبل 6 معطيات.
        \item عدّل دالة \textenglish{CalculateSalary()} في \textenglish{Trainer} لتضرب الراتب بـ 1.10.
        \item \textbf{كرر جميع هذه الخطوات الثلاث تماماً داخل فئة \textenglish{Manager}.}
    \end{itemize}

    \item
    \textbf{فحص الكود (في \textenglish{Main}):} \\
    \begin{itemize}
        \item قم بتحديث الكائنات القديمة لتشمل البريد الإلكتروني.
        \item تأكد من أن الراتب الجديد يُحسب بشكل صحيح بعد إضافة العلاوة.
    \end{itemize}
\end{enumerate}

\clearpage
\section*{المرحلة الرابعة: الوراثة كحل هندسي}
لحل مشكلة تكرار الكود وصعوبة التعديل، سنستخدم مفهوم الوراثة (Inheritance) لتمثيل علاقة التصنيف الهرمية.

\begin{enumerate}[itemsep=3em]
    \item
    \textbf{بناء فئة أساسية (\textenglish{Employee}):}
    \begin{itemize}
        \item انقل جميع الخصائص المشتركة إليها: \textenglish{Id}، \textenglish{FullName}، \textenglish{PhoneNumber}، \textenglish{EmailAddress}، \textenglish{BasicSalary}.
        \item اكتب دالة بناء واحدة هنا تهيئ هذه الخصائص.
        \item اكتب دالة \textenglish{CalculateSalary()} ودالة \textenglish{PrintInfo()} مرة واحدة فقط هنا.
    \end{itemize}

    \item
    \textbf{تحديث فئة المدرب (\textenglish{Trainer}):}
    \begin{itemize}
        \item اجعلها ترث من الفئة الأساسية (الفئة العليا) باستخدام النقطتين الرأسيتين \textenglish{:} للدلالة على \textenglish{Employee}.
        \item أبقِ خاصية \textenglish{Specialty} فقط (كتوسيع للفئة الفرعية).
        \item استخدم الكلمة المحجوزة \textenglish{base} في دالة البناء لتمرير المعطيات المشتركة لدالة بناء الفئة الأساسية.
    \end{itemize}

    \item
    \textbf{تحديث فئة المدير (\textenglish{Manager}):}
    \begin{itemize}
        \item اجعلها ترث من الفئة الأساسية \textenglish{Employee}.
        \item أبقِ خاصية \textenglish{Department} فقط.
        \item استخدم الكلمة المحجوزة \textenglish{base} في دالة البناء.
    \end{itemize}

    \item
    \textbf{فحص الكود (في \textenglish{Main}):}
    \begin{itemize}
        \item قم بتشغيل نفس الكود الموجود في الـ \textenglish{Main} من المرحلة الثالثة.
        \item يجب أن يعمل البرنامج بكفاءة، مع ملاحظة أن حجم الكود في فئتي \textenglish{Trainer} و \textenglish{Manager} أصبح قصيراً جداً وخالياً من التكرار بفضل إعادة استخدام الكود.
    \end{itemize}
\end{enumerate}

\vspace{1cm}
\begin{flushleft}
أرجو لكم وقتًا ممتعًا.

الأستاذ محمود اغبارية.
\end{flushleft}

\end{document}
